\section{Evaluation}
\label{sec:evaluation}

This section evaluates the proposed H3-based spatial indexing pipeline against a
rectangular grid baseline across two complementary dimensions: formal spatial
quality metrics and business-oriented performance indicators.
All results are produced by the \texttt{mart\_spatial\_comparison} dbt model running
on a synthetic dataset of 100{,}000 deliveries across 11 Ho Chi Minh City districts,
with four districts designated as structurally complex zones.

% ── 5.1 Experimental Setup ──────────────────────────────────────────────────

\subsection{Experimental Setup}

\paragraph{Data generation.}
The synthetic dataset is generated by \texttt{data\_generator/generate\_data.py}.
Each delivery location is drawn from a district centre with a Gaussian jitter of
$\pm 0.04°$ (${\approx}4.4$\,km radius), producing a realistic intra-district
spread.
Per-cell delay probability is derived from a hash of the GPS coordinate rather
than a uniform district flag, ensuring adjacent cells within the same district
exhibit heterogeneous delay profiles—mirroring real-world conditions where a
single block of narrow \textit{hẻm} alleys may be significantly more congested
than the surrounding streets.

\paragraph{Complex zones.}
Four districts (Quận 3, Quận 5, Quận 10, Gò Vấp) are designated as
\emph{spatially complex}.
Within these districts, individual cell delay probabilities range from 35\% to
85\% (mean $\approx 60\%$); within normal districts they range from 0\% to 22\%
(mean $\approx 11\%$).
When a delay occurs the additional latency is drawn uniformly from
$[15, 60]$\,min.
Delays are measured \emph{per stop}, independently of other stops on the same
route, so that the spatial signal is not confounded by cumulative driver drift.

\paragraph{Spatial methods compared.}
\begin{itemize}
  \item \textbf{H3 (resolution 9).} Uber's hierarchical hexagonal grid.
        Each cell covers ${\approx}0.1$\,km² with uniform area and
        equidistant neighbours by construction.
        Cell indices are computed via a Python UDF (\texttt{h3.latlng\_to\_cell})
        registered in DuckDB at load time.
  \item \textbf{Grid (0.01°).} Rectangular grid defined by
        $\lfloor\mathrm{lat}/0.01\rfloor$ and
        $\lfloor\mathrm{lon}/0.01\rfloor$.
        Each cell covers ${\approx}1.1$\,km² and is computed entirely in SQL
        using \texttt{FLOOR} arithmetic.
\end{itemize}

% ── 5.2 Spatial Quality Metrics ─────────────────────────────────────────────

\subsection{Spatial Quality Metrics}

Three formal metrics quantify how faithfully each method represents the
underlying continuous delivery geography.

\paragraph{Mean Distance Quantisation Error (MDQE).}
The average Haversine distance from each delivery's GPS coordinate to its
assigned cell centroid.
A lower MDQE indicates less information loss from snapping a continuous
coordinate to a discrete cell.

\paragraph{Spatial Distortion Coefficient (SDC).}
The standard deviation of Euclidean distances between a cell centre and its
immediate neighbours.
For H3, all neighbours are equidistant by construction, so $\text{SDC} = 0$.
For the rectangular grid, orthogonal neighbours are ${\approx}1.11$\,km apart
while diagonal neighbours are ${\approx}1.56$\,km apart, yielding a non-zero
SDC that biases proximity queries.

\paragraph{Computational Latency (CL).}
Microseconds per row, measured on the full 100{,}000-row dataset.
The H3 UDF incurs a Python-to-DuckDB inter-process call per row; the grid
method is pure in-engine SQL arithmetic.

\begin{table}[h]
\centering
\caption{Spatial quality metrics: H3 (res=9) vs.\ rectangular grid (0.01°).}
\label{tab:spatial_metrics}
\begin{tabular}{lrrr}
\toprule
\textbf{Method} & \textbf{MDQE (km)} & \textbf{SDC (km)} & \textbf{CL ($\mu$s/row)} \\
\midrule
H3 (res=9, ${\approx}0.1$\,km²)    & 0.128 & 0.000 & 3.31 \\
Grid (0.01°, ${\approx}1.1$\,km²)  & 0.422 & 0.228 & 0.07 \\
\midrule
\textit{H3 advantage}               & $3.3\times$ lower & equidistant & $49\times$ slower \\
\bottomrule
\end{tabular}
\end{table}

H3's MDQE of 0.128\,km versus the grid's 0.422\,km confirms that hexagonal
cells retain significantly more spatial precision: delivery coordinates are
snapped to a centroid that is, on average, 3.3 times closer to their true
position.
The zero SDC of H3 means that all neighbourhood relationships are
geometrically unbiased—a property that matters when the goal is to
identify the \emph{precise boundary} of a congested zone rather than
merely the district it falls within.

The grid's 49$\times$ lower computational latency (0.07 vs.\ 3.31\,$\mu$s/row)
is its main practical advantage.
At 100{,}000 rows the absolute difference is 0.33\,s—acceptable for a
nightly batch pipeline but worth noting for real-time ingestion scenarios.

% ── 5.3 Business Performance Metrics ────────────────────────────────────────

\subsection{Business Performance Metrics}

\paragraph{Hotspot detection.}
A cell is flagged as a \emph{hotspot} if its average delivery delay exceeds
15\,minutes.
Two retrieval metrics are computed across all hotspot cells:

\begin{itemize}
  \item \textbf{Capture rate} (\emph{recall}): percentage of all delayed
        deliveries (city-wide) that fall inside flagged cells.
  \item \textbf{Precision}: percentage of deliveries inside flagged cells
        that are themselves delayed.
\end{itemize}

\begin{table}[h]
\centering
\caption{Business performance metrics: H3 vs.\ grid.}
\label{tab:business_metrics}
\begin{tabular}{lrrrrr}
\toprule
\textbf{Method} & \textbf{Cells} & \textbf{Hotspot cells}
  & \textbf{Capture rate (\%)} & \textbf{Precision (\%)}
  & \textbf{Within-cell delay $\sigma$ (min)} \\
\midrule
H3 (res=9)  & 3{,}550 & 1{,}000 & 60.3 & 48.4 & 13.67 \\
Grid (0.01°) &   393  &   105  & 61.5 & 47.1 & 13.98 \\
\bottomrule
\end{tabular}
\end{table}

Both methods achieve similar macro recall (${\approx}60\%$) and precision
(${\approx}48\%$), confirming that the rectangular grid is not categorically
inferior for coarse hotspot detection.
The key differentiator emerges when considering granularity: H3 identifies
1{,}000 distinct problem zones across 3{,}550 mapped cells, whereas the grid
collapses the same geography into only 393 cells of which 105 are flagged.
A single flagged grid cell covers an area roughly eleven times larger than an
H3 cell, making it impossible to distinguish which \emph{part} of that cell
drives the delay—the information that operations teams actually need to re-route
drivers or apply zone-specific surcharges.

\paragraph{Within-cell delay homogeneity.}
The average within-cell standard deviation of delay minutes (13.67 for H3
vs.\ 13.98 for the grid) is marginally lower for H3.
A lower standard deviation indicates that cells are more internally
homogeneous—deliveries within the same H3 cell experience similar delays,
supporting the assumption that the cell captures a coherent congestion zone.

\paragraph{Margin gap concentration.}
Both methods concentrate a material share of the total margin gap inside
flagged cells (31.6\% for H3, 33.4\% for the grid), demonstrating that
spatial indexing—regardless of shape—successfully identifies areas where
operational losses are disproportionately high.
Because H3 cells are smaller, the flagged zones represent a more surgically
targeted set of addresses for pricing interventions.

% ── 5.4 Discussion ──────────────────────────────────────────────────────────

\subsection{Discussion}

The evaluation reveals a clear trade-off between geometric fidelity and
computational cost.

\textbf{H3 is the better choice for spatial analytics} when the objective is
to identify and act on sub-kilometre congestion patterns.
Its 3.3$\times$ MDQE advantage and zero SDC mean that zone boundaries are
determined by the data rather than by the arbitrary alignment of a
degree-based grid.
The 1{,}000 distinct hotspot cells it produces allow operations to attach
zone-specific rules (surcharges, driver briefings, routing constraints) at
a level of precision that a 105-cell grid report cannot support.

\textbf{The grid retains a role} in latency-sensitive pipelines or
resource-constrained environments where the H3 Python UDF cannot be
registered.
Its pure-SQL implementation is portable across any SQL engine without
extension dependencies, and its 49$\times$ speed advantage becomes meaningful
at ingestion scales above ${\sim}10^7$ rows per batch.

\textbf{Limitations of the evaluation.}
The dataset is fully synthetic: district boundaries are approximated by
Gaussian jitter around a single centre point, and delay probabilities are
assigned by a coordinate hash rather than a real congestion model.
Real HCMC delivery data would exhibit spatial autocorrelation, time-of-day
effects, and irregular district shapes that may shift the relative
performance of the two methods.
Future work should validate the pipeline on live GPS traces and consider
adaptive H3 resolution—using resolution 8 (${\approx}0.7$\,km²) in sparse
suburban areas and resolution 10 (${\approx}0.015$\,km²) in dense urban cores.

% ── 5.5 Summary ─────────────────────────────────────────────────────────────

\subsection{Summary}

Table~\ref{tab:summary} consolidates the evaluation findings across all
dimensions.

\begin{table}[h]
\centering
\caption{Consolidated evaluation summary.}
\label{tab:summary}
\begin{tabular}{lcc}
\toprule
\textbf{Criterion} & \textbf{H3 (res=9)} & \textbf{Grid (0.01°)} \\
\midrule
Cell area              & ${\approx}0.1$\,km²   & ${\approx}1.1$\,km²  \\
MDQE                   & 0.128\,km             & 0.422\,km            \\
Geometric uniformity   & \checkmark (SDC=0)    & $\times$ (SDC=0.228) \\
Hotspot capture rate   & 60.3\%                & 61.5\%               \\
Hotspot precision      & 48.4\%                & 47.1\%               \\
Actionable granularity & 1{,}000 zones         & 105 zones            \\
Computational latency  & 3.31\,$\mu$s/row      & 0.07\,$\mu$s/row     \\
SQL-only implementation & $\times$             & \checkmark           \\
\bottomrule
\end{tabular}
\end{table}

H3 at resolution 9 is recommended as the primary spatial index for the
HCMC last-mile analytics pipeline.
It provides the granularity required for actionable operational decisions
while maintaining acceptable ingestion latency for a nightly batch workload.
The grid baseline remains a useful fallback and a meaningful benchmark for
validating that spatial indexing—in any form—outperforms purely district-level
reporting.
